\providecommand{\todo}[1]{\textbf{[TODO: #1]}}

\section{ความเป็นมาและความสำคัญของโครงงาน}

การตรวจชิ้นเนื้อด้วยตาเปล่า (Gross Examination) เป็นขั้นตอนแรกของการตรวจสิ่งส่งตรวจทางศัลยพยาธิวิทยา ในกรณีมะเร็งเต้านม พยาธิแพทย์หรือผู้ช่วยพยาธิแพทย์ต้องวัดขนาดก้อนเนื้อในสามมิติ ประเมินระยะห่างของก้อนจากขอบตัด (Surgical Margin) ทั้ง 6 ด้านในชิ้นเนื้อแบบ Lumpectomy และตรวจนับจำนวนพร้อมทั้งวัดขนาดต่อมน้ำเหลือง ข้อมูลเหล่านี้เมื่อผสานกับการตรวจทางจุลพยาธิวิทยาจะถูกบันทึกตามโปรโตคอลมาตรฐานของ College of American Pathologists (CAP) เพื่อใช้กำหนดระยะของโรคตามระบบ pTNM \cite{cap_breast,ajcc2017}

ในทางปฏิบัติ ผู้ตรวจต้องสวมถุงมือที่สัมผัสสิ่งส่งตรวจ สารคัดหลั่ง และฟอร์มาลิน จึงไม่สามารถพิมพ์บันทึกข้อมูลผ่านแป้นพิมพ์คอมพิวเตอร์ได้โดยสะดวก แนวทางเดิมมักใช้การบันทึกเสียงหรือจดบันทึกลงกระดาษแล้วพิมพ์ลงระบบในภายหลัง ซึ่งก่อให้เกิดความล่าช้าและเพิ่มโอกาสเกิดความผิดพลาดในการถ่ายโอนข้อมูล (Transcription Error) สำหรับข้อมูลวิกฤต เช่น ขนาดก้อนและระยะห่างจากขอบตัด ขณะที่ระบบรู้จำเสียงพูดทางคลินิกทั่วไปยังพบอัตราความผิดพลาดเฉลี่ยร้อยละ 7.4 ในผลดิบก่อนการตรวจทาน \cite{zhou2018errors}

นอกจากนี้ ระบบแปลงเสียงเป็นข้อความส่วนใหญ่ประมวลผลบนคลาวด์ ซึ่งต้องส่งสัญญาณเสียงที่มีข้อมูลสุขภาพของผู้ป่วยออกนอกเครือข่ายภายใน จึงมีความเสี่ยงต่อข้อกำหนดด้านการคุ้มครองข้อมูลส่วนบุคคลตามพระราชบัญญัติคุ้มครองข้อมูลส่วนบุคคล พ.ศ. 2562 (PDPA) และ HIPAA \cite{pdpa2562,hipaa} ประกอบกับสภาพแวดล้อมจริงในห้องตรวจชิ้นเนื้อมีเสียงรบกวนต่อเนื่องจากพัดลมของตู้ดูดควันไอระเหยฟอร์มาลิน (Fume Hood) และลักษณะการพูดของแพทย์มีความเร็ว มีการพูดแก้ไขตัวเลขหรือมิติกลางประโยค แม้แบบจำลอง Whisper \cite{radford2023whisper} จะทำงานแบบออฟไลน์และทนต่อสัญญาณรบกวนได้ดี แต่เมื่อประมวลผลบนหน่วยประมวลผลกลาง (CPU) ทั่วไปจะใช้เวลาแฝงสูง และมีความเสี่ยงต่อการสร้างข้อความที่ไม่มีอยู่ในสัญญาณเสียงจริง (Hallucination) ราวร้อยละ 1.4 ของช่วงเสียง \cite{koenecke2024careless}

โครงงานนี้จึงพัฒนาระบบ PathoWhisper เพื่อทำหน้าที่เป็นเครื่องมือสนับสนุนการบันทึกข้อมูลทางพยาธิวิทยาแบบออฟไลน์บนเครื่องคอมพิวเตอร์ส่วนบุคคล โดยผสานการกรองสัญญาณรบกวนจากตู้ดูดควัน การถอดความด้วย Faster-Whisper ที่ผ่านการทำ INT8 Quantization ร่วมกับการกำหนดบริบทคลังคำศัพท์ CAP และการสกัดข้อมูลลงแบบฟอร์มด้วยกฎกำหนดแน่นอน (Deterministic Extraction) ก่อนส่งออกรายงานสรุปในรูปแบบเอกสาร PDF

\section{วัตถุประสงค์ของโครงงาน}

\begin{enumerate}
    \item เพื่อพัฒนาระบบแปลงเสียงบรรยายผลการตรวจชิ้นเนื้อทางพยาธิวิทยาเป็นข้อความ ที่สามารถประมวลผลแบบออฟไลน์บนเครื่องคอมพิวเตอร์ส่วนบุคคล
    \item เพื่อบูรณาการกระบวนการกรองสัญญาณรบกวนจากตู้ดูดควันเข้ากับการปรับแต่งการถอดความ (PathoWhisper INT8) ให้เหมาะสมกับศัพท์เฉพาะทางศัลยพยาธิวิทยา
    \item เพื่อพัฒนาตัวสกัดข้อมูลทางคลินิกแบบกำหนดแน่นอนที่แปลงข้อความถอดความลงในแบบฟอร์มสำหรับมะเร็งเต้านม และสร้างเอกสารรายงาน PDF โดยอัตโนมัติ
    \item เพื่อประเมินประสิทธิภาพของระบบเทียบกับแบบจำลองพื้นฐาน ทั้งในมิติอัตราความผิดพลาดระดับคำ (WER) เวลาในการประมวลผล และความถูกต้องของการสกัดข้อมูลลงฟิลด์
\end{enumerate}

\section{ขอบเขตของโครงงาน}

\subsection{ด้านข้อมูล}
\begin{itemize}
    \item ชุดข้อมูลทดสอบ 1,000 กรณีศึกษา แบ่งเป็น 10 หมวดหมู่ หมวดละ 100 กรณี
    \item ข้อมูลเสียงมีความยาวเฉลี่ย 21.48 วินาทีต่อกรณี ครอบคลุมรูปแบบการรายงานปกติ การรายงานสลับลำดับ การแก้ไขข้อมูลกลางประโยค สภาวะเสียงรบกวนจากตู้ดูดควันที่ระดับ SNR $-10$ dB และการพูดเร็วเกิน 200 คำต่อนาที
    \item สัญญาณเสียงสังเคราะห์จากสคริปต์ข้อความตามโครงสร้างโปรโตคอล CAP ด้วยระบบสังเคราะห์เสียงภาษาอังกฤษ (Microsoft Edge TTS: สำเนียง en-US สองรูปแบบเสียง) ร่วมกับการซ้อนทับสัญญาณเสียงรบกวนที่บันทึกจากตู้ดูดควัน
    \item การทดสอบใช้ข้อมูลและค่าตัวเลขสมมุติตามเกณฑ์มาตรฐาน CAP โดยไม่มีการเก็บรวบรวมข้อมูลส่วนบุคคลหรือสิ่งส่งตรวจจริงจากผู้ป่วย จึงไม่ต้องผ่านกระบวนการรับรองจริยธรรมการวิจัยในมนุษย์ในขั้นตอนนี้
\end{itemize}

\subsection{ด้านเทคนิค}
\begin{itemize}
    \item แบบจำลองการถอดความ: Faster-Whisper รุ่น Small บนเฟรมเวิร์ก CTranslate2 ประมวลผลแบบ INT8 บน CPU โดยปรับแต่งพฤติกรรมผ่านการกำหนดข้อความนำบริบท (Prompt Conditioning) ตามคลังศัพท์ CAP โดยไม่มีการปรับแก้ค่าน้ำหนัก (Fine-tuning)
    \item การประมวลผลเสียง: การลดสัญญาณรบกวนสเปกตรัมด้วยตัวกรอง afftdn และการตัดช่วงเงียบสัญญาณเสียงด้วย silenceremove ผ่าน FFmpeg
    \item การสกัดข้อมูล: ใช้นิพจน์ปรกติ (Regular Expression) และการวิเคราะห์บริบทแบบกำหนดแน่นอน ครอบคลุมโครงสร้างรายงานมะเร็งเต้านมตามเกณฑ์ CAP จำนวน 15 ฟิลด์
    \item สภาพแวดล้อมฮาร์ดแวร์: ประมวลผลบนหน่วยประมวลผลกลาง (CPU) 8 เธรด โดยไม่ใช้หน่วยประมวลผลกราฟิก (GPU) ภายใต้หน่วยความจำหลัก 16 GB
\end{itemize}

\subsection{ด้านระบบ}
\begin{itemize} 
\item เว็บแอปพลิเคชันติดตั้งและให้บริการภายในเครื่องของผู้ใช้ (Locally Hosted Application) ผ่านโพรโทคอล HTTPS
\item รองรับการบันทึกเสียงสดผ่านไมโครโฟนบนเบราว์เซอร์และการประมวลผลไฟล์เสียงที่บันทึกไว้ล่วงหน้า
\item ระบบส่งออกผลลัพธ์เป็นข้อความถอดความดิบ ข้อมูลที่สกัดลงแบบฟอร์ม 15 ฟิลด์ และรายงานสรุปในรูปแบบเอกสาร PDF 
\item ระบบจัดเก็บข้อมูลภายในเครื่องโดยใช้ PostgreSQL เป็นฐานข้อมูลหลัก และมี SQLite เป็นฐานข้อมูลสำรองในกรณีออฟไลน์
\end{itemize}

\subsection{ข้อจำกัดของโครงงาน}
\begin{itemize}
    \item ครอบคลุมเฉพาะการตรวจชิ้นเนื้อมะเร็งเต้านมด้วยตาเปล่าในหัตถการ Mastectomy และ Lumpectomy
    \item รองรับเฉพาะการบรรยายเป็นภาษาอังกฤษทางการแพทย์
    \item การประเมินระยะห่างขอบตัดครอบคลุมเฉพาะด้านลึก (Deep Margin)
    \item ผลการประเมินอิงกับชุดข้อมูลเสียงสังเคราะห์ในสภาวะจำลอง ซึ่งทำหน้าที่เป็นเครื่องมือสนับสนุนการลงข้อมูลเบื้องต้น โดยแพทย์ผู้ตรวจต้องตรวจสอบความถูกต้องก่อนการยืนยันผลเสมอ
\end{itemize}